\documentclass[final,5p,times,twocolumn]{elsarticle}
% This is the two-column Procedia CIRP submission layout.
% For a wider single-column draft to mark up, swap the line above for:
%   \documentclass[preprint,times]{elsarticle}

\usepackage{amsmath,amssymb}
\usepackage{graphicx}
\usepackage{booktabs}
\usepackage{siunitx}
\usepackage{lineno}
\usepackage{array}
\usepackage{tabularx}
\usepackage{rotating}
\usepackage{xurl}
\usepackage{float}
\usepackage{placeins}
\usepackage{tikz}
\usetikzlibrary{arrows.meta,positioning,shapes.geometric,fit}
\usepackage[T1]{fontenc}
\usepackage{xcolor}
\usepackage{hyperref}
\hypersetup{colorlinks=true,linkcolor=blue,citecolor=blue,urlcolor=blue}

\newcolumntype{M}{>{\hsize=1.5\hsize\raggedright\arraybackslash}X}
\newcolumntype{D}{>{\hsize=0.6\hsize\centering\arraybackslash}X}
\newcolumntype{T}{>{\hsize=0.6\hsize\centering\arraybackslash}X}
\newcolumntype{S}{>{\hsize=1.5\hsize\raggedright\arraybackslash}X}
\newcolumntype{F}{>{\hsize=0.8\hsize\raggedright\arraybackslash}X}

\journal{[Journal name -- see cover letter]}

% compact reference list, as in the published Procedia page
\setlength{\bibsep}{0pt plus 0.3ex}

\begin{document}

\begin{frontmatter}

\title{Which process faults can be told apart? Diagnosability limits of
in-situ monitoring in laser powder bed fusion}

%% ---------------------------------------------------------------------------
%% TODO before submission: complete author list, affiliation and corresponding
%% author details. The entry below is a placeholder.
%% ---------------------------------------------------------------------------
\author[a]{Soham Jain\corref{cor1}}
\ead{f20240343@pilani.bits-pilani.ac.in}
\cortext[cor1]{Corresponding author.}
\affiliation[a]{organization={[Affiliation to be completed]},
                city={[City]},
                country={[Country]}}

\begin{abstract}
In-situ monitoring is standard on metal additive manufacturing systems, and
classifiers report high accuracies for detecting defects from photodiode and
imaging signals. Those results answer a question about outcomes: which defect
formed. Specifying hardware requires an answer one step upstream: which process
\emph{causes} a configuration can separate in principle, and which remain
confounded whatever classifier is applied. Where two causes produce signatures
differing by less than the sensor noise, no volume of training data distinguishes
them: the limit is structural, not algorithmic. This paper applies structural and
quantitative diagnosability analysis, mature in automotive and aerospace process
control, to laser powder bed fusion. A calibrated reduced-order thermal model
gives the response of a coaxial suite to five injected faults (laser power and
scan speed deviation, thin powder layers, degraded shielding gas flow and
geometry-driven heat accumulation) with separability measured as an angle
between noise-whitened signatures and each sensor's relative noise derived from
its transduction chain rather than assumed. Laser power and gas degradation prove
\emph{exactly} confounded instantaneously, by algebra rather than noise, and
seven candidate observables leave the attained rank at three against a structural
bound of four. Regressing the same three sensors onto known build context raises
the rank to five and the worst pair from $0.00^\circ$ to $24.90^\circ$ at no
hardware cost, outperforming two idealised direct measurements of the causes
themselves. Matrices are reported for five configuration classes found on current
commercial machines.
\end{abstract}

\begin{keyword}
Additive manufacturing \sep Process monitoring \sep Diagnosability \sep
Fault diagnosis \sep Laser powder bed fusion
\end{keyword}

\end{frontmatter}

%% ===========================================================================
\section{Introduction}

Laser powder bed fusion (LPBF) builds a part by melting thousands of short track
segments per layer, and a modest excursion in any of several independent inputs
can produce a rejected part. Manufacturers have responded by fitting in-situ
monitoring as standard, with coaxial photodiodes, melt pool cameras, pyrometers,
increasingly acoustic and powder-bed imaging, and a large literature maps
those signals onto defects with high reported accuracy.

Almost all of that work answers a question about \emph{outcomes}: given the
signals, which defect formed? That matters for part acceptance, but not for the
machine builder specifying hardware or the process engineer facing a drifting
build. Those concern \emph{causes}, that the laser is under-delivering, the
recoater short-feeding, the gas flow degraded, and which it is determines what
gets corrected. The distinction is not semantic, because causes can be mutually
indistinguishable in a way outcomes are not: if two of them drive the readings
along the same direction in signal space, the readings carry no information
separating them and no classifier, however trained, recovers the difference. That
limit is structural, a property of the measurement and the physics rather than of
the algorithm, where a pair separable in principle but misclassified in practice
is an algorithm that can be improved.

Diagnosability analysis is the established discipline for deciding which case
holds. Developed within model-based fault detection and isolation, it asks,
before any estimator is built, which faults a measurement set can detect and
which it can isolate from one another \cite{Krysander2008}. The methods
are standard where sensor sets are costly and must be justified in advance, but
to the authors' knowledge have not been applied to LPBF.

This paper applies them. A calibrated reduced-order thermal model is driven by
five injected fault modes and the signatures compared pairwise against noise,
yielding a diagnosability matrix per configuration that divides fault pairs into
those resolvable from existing signals, those needing an added measurement and
those structurally irreducible. The contributions are the first such matrix for
the process, to the authors' knowledge; noise figures derived from each sensor's
transduction chain rather than assumed; a demonstration that laser power and gas
degradation are \emph{exactly} confounded instantaneously; the finding that the
binding limit is the melt pool rather than the instrumentation; and that build
structure already present dissolves it at no hardware cost.

\section{Related work}

\subsection{What in-situ monitoring classifies}

Melt pool monitoring with coaxial photodiodes and cameras is a most widely
deployed configuration for flaw classification in LPBF processes.
Scime and Beuth identified signatures indicative of keyholing porosity and
balling from high-speed visible imaging \cite{Scime2019}. Analysis also exists across modalities and alloys, including mechanical-property prediction from
coaxial signals \cite{Raj2023}. The predicted quantity is in each case an outcome
or a material state, and whether two distinct excursions could have produced the
same signal is never posed.

Experimental methods have been applied to LPBF to get the relationship between process inputs and obtained outputs. RAISE-LPBF reconstructs laser power and scan speed independently from
on-axis video on 316L, sampling both on every scan line from continuous distributions
so interactions can be studied \cite{Blanc2024}. Reijonen reports coaxial
photodiode response to power, speed, focal position, scan orientation and
shielding gas flow, finding power approximately linear, speed a power law and the
rest non-linear and interactive \cite{Reijonen2024}. Both characterise
\emph{sensitivity}; identifiability, meaning whether sensitivities to different causes
are linearly distinguishable, is what neither addresses.

A third class observes the powder bed, attacking causes the melt pool reports
only indirectly \cite{Liu2022,Bugatti2022}, and returns in
Section~\ref{sec:added} as a target rather than a modelled option.

\subsection{Use of lumped process metrics}

\emph{Lumped} process quantities have are increasingly being used in place of component values. Volumetric energy density (VED) is a common choice, and
state-of-the-art parameter estimation from melt pool imagery predict
energy density \emph{class} rather than power and speed separately \cite{kNN2025}.

While analytically helpful, lumped quantities often discard information. Guo et al.\ held input
energy density constant while varying power and speed together and saw melt pool
dimensions change by orders of magnitude with transitions between melting
regimes \cite{Guo2019}. The 2022 Additive Manufacturing Benchmark found
single-track depth correlated with VED while width did not \cite{Weaver2024}.
Both are reproduced in Section~\ref{sec:valid}. If the signals cannot distinguish power from speed, then only some combination
of the two is estimable. Energy density, though, is a lumping adopted by convention rather than derived from what any particular suite measures. Section~\ref{sec:valid} finds melt pool depth only partly a function of it and width not at all. Whether the components are in fact separable for a given
sensor set, is the diagnosability question we investigate in this work.

\subsection{Diagnosability analysis}

Structural analysis uses only which variables appear in which equations to
decide which faults are detectable and which separable \cite{Krysander2008}. It
is immune to parameter error but optimistic: structurally separable faults can
still be numerically indistinguishable. Quantitative analysis settles that from
the geometry of the linearised fault-to-sensor map. Both are run here.

%% ===========================================================================
\section{Methodology}
%% ===========================================================================

\subsection{Forward model and calibration}

Melt pool geometry is computed from the Eagar--Tsai solution for a moving Gaussian
heat source on a semi-infinite solid, extended with the Gladush--Smurov keyhole
depth correction following Honarmandi et al.\ \cite{Hon2021}. Four parameters: effective absorptivity, effective conductivity, heat capacity, the keyhole
constant are fitted to the seven IN718 bare-plate single-track cases of the
2022 Additive Manufacturing Benchmark \cite{Weaver2024}, giving \textbf{11.7\,\%
mean error on melt pool width and 13.1\,\% on depth}. All results below are evaluated at that benchmark's own baseline condition,
285\,W, 960\,mm\,s$^{-1}$ and a 67\,\textmu m D4$\sigma$ spot, where the model under-predicts width by 15.5\,\% and over-predicts depth by
1.7\,\%. The keyhole correction
is necessary as all seven cases are keyhole with depth-to-width
ratios 0.72--2.14 against the criterion's 0.67, and pure conduction gives 23.6\,\% depth error. Equation~\eqref{eq:jac} is evaluated numerically, one fault at a time:
%
\begin{equation}
\frac{\partial y_s}{\partial f_k} \;\simeq\;
\frac{y_s(f_k = +h) - y_s(f_k = -h)}{2h} + O(h^{2}).
\label{eq:cd}
\end{equation}
%
Too large a step measures curvature rather than slope, so $h$ is chosen by
sweeping it. At $h=0.02$ the entries still shift 39.3\,\% per halving and below
$h=0.01$ they settle, and the shift then falls fourfold per halving, as the
$O(h^{2})$ term requires. All results use $h=0.002$, where halving again moves
them by 0.007\,\%.

\subsection{Fault definitions}

Five faults are modelled, each defined as a \emph{physical} quantity acting on a
model input. For
magnitude $f$: \emph{power} scales laser power, $P(1+f)$; \emph{speed} scales
scan speed, $v(1+f)$; \emph{gas} is a plume optical depth attenuating effective
absorptivity, $A\,e^{-f}$; \emph{geometry} is a preheat rise in units of 100\,K,
$T_0+100f$; \emph{layer} is a fractional layer-thickness deviation, positive
denoting a \emph{thinner} layer as produced by recoater short-feed, acting at
once on absorptivity, effective conductivity and effective power,
$A\,e^{c_A f}$, $\kappa\,e^{c_\kappa f}$ and $P\,e^{c_H f}$.

The absorptivity
coupling, $-0.230$, follows from the powder-to-solid absorptivity ratio of
approximately two measured calorimetrically at 1070\,nm \cite{Trapp2017} with a
radiative penetration depth of a few particle diameters, implying a bare-solid
absorptivity of 0.258 consistent with the fitted value. The powder-melting
enthalpy coupling, $+0.186$, follows from an energy balance in which melting the
layer consumes 15.7\,\% of the absorbed energy, so a thinner layer frees
proportionally more for penetration. The conductivity coupling remains an unmeasured estimate, swept across its range
in Section~\ref{sec:robust} with every accumulated-regime result invariant to it,
but is important as the other two channels nearly cancel on
the absorbed-power product ($-0.230+0.186=-0.044$), so without it the layer fault
acts only through that product and joins the power/gas degeneracy
\emph{exactly} as its Jacobian column becomes $-0.044$ times the power column on
all three sensors, and power/layer and layer/gas both read $0.00^\circ$. This one
unmeasured coupling is all that distinguishes a thin layer from a weak laser at
an instant. All couplings are applied in exponential form $A\,e^{cf}$ as they may agree with
$A(1+cf)$ to first order but are unable to pass through zero, where the linear form
would drive thermal diffusivity negative and the integrand imaginary.

\subsection{Sensor models}

Three coaxial data points are computed from the temperature field: a
band-integrated photodiode (700--1000\,nm), melt pool area, and a two-colour
pyrometer (750/950\,nm) evaluated as the ratio of integrated band powers and
inverted for temperature. Spectral radiance follows Planck's law, temperature
being capped at the boiling point before any radiance is computed. Four
additional options are defined for Section~\ref{sec:melt}: the camera's width and length
reported separately, a short-wave infrared photodiode and a wide-band
pyrometer.

\subsection{Diagnosability measures}
\label{sec:measures}

The relative-sensitivity Jacobian is
%
\begin{equation}
J_{sk} \;=\; \frac{1}{y_s}\,\frac{\partial y_s}{\partial f_k},
\label{eq:jac}
\end{equation}
%
so entry $(s,k)$ is the fractional change in sensor $s$ per unit of fault $k$.
Each row is then
divided by that sensor's relative noise $\sigma_s$, derived in
Section~\ref{sec:noise}, to give the whitened matrix $J^{\mathrm{w}}$. This makes the rows be in a common unit and angles between columns physically meaningful.

Three quantities are reported throughout. The \emph{angle} between whitened columns,
\begin{equation}
\theta_{kl} = \arccos\frac{|\langle J^{\mathrm{w}}_{:k},
J^{\mathrm{w}}_{:l}\rangle|}{\|J^{\mathrm{w}}_{:k}\|\,\|J^{\mathrm{w}}_{:l}\|},
\label{eq:angle}
\end{equation}
measures separability: at $0^\circ$ two faults move the readings along one line and no estimator separates them, at $90^\circ$ they are orthogonal. A
$10^\circ$ floor marks a pair resolved, a convention whose influence is reported. The \emph{rank} of $J^{\mathrm{w}}$ counts how many fault directions are seen at all, a singular value below $10^{-6}$ of the largest counting as none. Rank alone misleads, scoring a direction barely above that cut equal to the strongest, so the \emph{condition number} of the retained singular values,
\begin{equation}
\kappa = \sigma_1/\sigma_r,
\qquad \sigma_1 \ge \dots \ge \sigma_r > 10^{-6}\sigma_1,
\label{eq:cond}
\end{equation}
is reported beside it: the weakest direction seen responds $\kappa$ times less than the strongest. Structural analysis runs in parallel on an eleven-equation representation using minimal structurally overdetermined sets, as an independent check.

\subsection{Build context and the enlarged observation}
\label{sec:context}

A build is not a single instant. It comprises thousands of segments whose
toolpath, layer index, gas-flow direction and local geometry are known to the
machine, and the five faults do not vary identically across them: laser
degradation drifts with build time, galvanometer error concentrates at turns, a
recoater short-feed recurs with layer parity, gas degradation follows the flow
axis, heat accumulation follows support depth. Treat the build rather than the
instant as the observation. Let $G\in\mathbb{R}^{N\times 5}$ hold a constant and
the last four of those quantities as known regressors over $N$ segments, and
$M\in\mathbb{R}^{N\times 5}$ the modulation profile of each fault. Regressing
each sensor's relative residual onto $G$ returns five coefficients per sensor in
place of one reading, so three sensors give 15 observables, with Jacobian
%
\begin{equation}
J^{\mathrm{big}}_{(s,j),k} = J_{sk}\,B_{jk},
\qquad B = (G^{\!\top}G)^{-1}G^{\!\top}M .
\label{eq:big}
\end{equation}
%
Coefficient $(s,j)$ carries noise $\sigma_{(s,j)} =
\sigma_s\sqrt{[(G^{\!\top}G)^{-1}]_{jj}}$, not $\sigma_s$. Were the regressors
orthogonal that factor would be $1/\sqrt{N}$, or 0.109 at $N=84$; correlation
among them lifts the worst to 0.136, so fitting a build buys a seven- to
ninefold noise reduction. Section~\ref{sec:robust} gives the dependence on $N$.

The profiles in $M$ are physical claims rather than measurements, and they are
\emph{machine-specific}: a different nozzle geometry, scan strategy or part
gives different patterns, so there is no universal set of five to go and find.
One can nonetheless be calibrated against published data, and that inversion is
the recipe by which a reader calibrates the rest on their own machine. The
AM Bench pad scans \cite{Weaver2024}, under a 4.3\,m\,s$^{-1}$ cross-flow, show
melt pool depth falling from 175 to 125\,\textmu m on tracks scanned parallel to
the gas flow rather than anti-parallel, a 27.2\,\% mean reduction against
standard deviations of 4--11\,\textmu m, with perpendicular pads as a control at
2.4--4.7\,\%. Weaver et al.\ attribute it to plume by-products attenuating,
scattering and lensing the beam; this model represents the attenuation alone.
The gas fault is a plume optical depth, transmitting $e^{-f}$, so each measured
depth change inverts to the severity reproducing it, through depth, where the
model errs by $+1.7$\,\% against $-15.5$\,\% on width. The control's $-3.6$\,\%
gives $f=0.046$, a plume taking 4.5\,\% of the beam; the parallel mean's
$-27.2$\,\% gives $0.404$, taking 33\,\%: an \textbf{8.8$\times$ swing in gas
severity inside one healthy build}, from scan direction alone. A profile of
\emph{contrast} $c$ varies as $e^{cz}$ with $z$ over $[-1,1]$ across the build,
so it swings $e^{2c}$ between mildest segment and worst; 8.8$\times$ is $c=1.08$
against the 0.90, or 6.05$\times$, assumed here, so \textbf{the one profile that
can be checked is conservative by 20\,\%}. Three qualifications: the driver is
scan direction rather than position along the flow, it is baseline plume
behaviour rather than a fault, and the inversion is model-mediated. The gas
signature also alternates at track scale where a recoater signature alternates
at layer scale, the distinctness Section~\ref{sec:robust} shows the method
needs.

\subsection{Sensor noise, derived from transduction}
\label{sec:noise}

The angles of Section~\ref{sec:measures} are angles between \emph{whitened}
columns, and whitening divides each sensor's row by its own relative noise. That
division is what puts a photodiode's radiometric power, a camera's square metres
and a pyrometer's kelvin into one common unit, so the noise figures are half of
the measurement geometry, not a detail applied after it. No published value for
these instruments' in-process relative noise was found; the one citable
figure, 2.3\,\%, is the precision of a post-mortem cross-section
(6.2\,\textmu m at $k=2$ on 136.3\,\textmu m \cite{Weaver2024}) rather than of a
signal acquired during a build. The three are therefore derived from a single
assumption, that every radiometric channel carries the same fractional noise
$\varepsilon$ whatever its physical cause, with the rest following from how each
instrument converts radiance into a reading. Only ratios enter the whitening, so
$\varepsilon$ cancels and no absolute noise level is claimed.

\textbf{Photodiode.} Its reading is the light in its band, summed, so a
fractional error in the light is the same fractional error in the reading:
$\sigma_{\mathrm{pd}} = \varepsilon$, the unit the other two are quoted in.

\textbf{Two-colour pyrometer.} It divides the power in one band by the power in
another and inverts the ratio for temperature. Two independent channel errors
put $\sqrt{2}\,\varepsilon$ on the ratio, and the ratio changes by 1.333\,\% for
each 1\,\% change in temperature, so the inferred temperature carries
$\sqrt{2}\,\varepsilon/1.333 = 1.061\,\varepsilon$: as a fraction of its own
reading, slightly noisier than the photodiode it is built from. The reason is
that 750 and 950\,nm are close together, either band alone changes by
5--6\,\% per 1\,\% of temperature, so their ratio keeps under a quarter of the
sensitivity available, and a single calibrated band would do five to six times
better. That is the price of the ratio's one virtue, needing no absolute
calibration.

\textbf{Camera.} It reports melt pool extent, found by thresholding at the
liquidus, where a fractional brightness error is indistinguishable from the
threshold sitting at a different temperature, $\delta \ln T = -\varepsilon /
(\partial \ln L/\partial \ln T)$. Planck's law is steep there, a 1\,\%
temperature change alters the radiance at 850\,nm by 10.52\,\%, so
$\varepsilon = 1\,\%$ displaces the apparent isotherm by only 0.095\,\% in
temperature. Carrying that through the model's own dimension solver moves width
by 0.030\,\% and length by 0.112\,\%, and the reported area is their product:
$\sigma_{\mathrm{cam}} = 0.030 + 0.112 = 0.142\,\varepsilon$. Thresholding on a
steep function localises an edge well, making the camera the quietest of the
three by a factor of seven and its width axis alone by a factor of thirty-four,
which is why Section~\ref{sec:added} finds the two axes worth reporting
separately.

The derived set is $1 : 0.142 : 1.061$. Two choices inside it were taken in the
direction that costs this analysis: the camera's extent is read from one
threshold crossing rather than averaged around the perimeter, and width and
length are treated as independent although a common brightness error moves both
together, which understates what separating them is worth. Both limits are swept
in Section~\ref{sec:robust}, alongside a case in which the dominant error is a
slow gain drift shared by the pyrometer's two channels, where the ratio then
cancels it and the pyrometer becomes the quietest of the three rather than the
noisiest. The worst pair over a build moves by under half a degree across all
four sets.

%% ===========================================================================
\section{Results}
%% ===========================================================================

\subsection{Validation}
\label{sec:valid}

This paper's quantities are angles between whitened Jacobian columns,
\emph{directions} in signal space, so what the model must get right is how
differently the faults move the readings, not how far. Guo et al.\ held
volumetric energy density constant while raising laser power and scan speed
together, and found the melt pool changing by orders of magnitude, passing
between melting regimes \cite{Guo2019}. The same sweep here, power 150 to
500\,W at compensated speed, changes in the same direction and does change
regime, but by a factor of a few rather than orders: depth moves 1.11$\times$
across the six conduction conditions and reaches 2.37$\times$, pool volume
5.64$\times$, only once the seventh keyholes. The model's whole variation is
that single discrete transition. Absorption varying with the geometry of the
vapour depression, which Guo et al.\ identify as the cause, is absent from a
conduction solution.

Its absence runs in a known direction. An absorption tracking the depression
responds to power and to speed unequally, a channel the two faults would not
share, and Section~\ref{sec:limits} measures what adding such a channel does to
the gas fault: it \emph{opens} a confounded pair. The model understates how
differently power and speed act, so the confounding reported here is an upper
bound. The same sweep also settles which dimension carries the information, used
in Section~\ref{sec:added}: width and depth move 1.11$\times$ each, in opposite
directions, while length varies 2.53$\times$. The cross-section is nearly blind
to how power and speed were traded off and the length is not, again a
statement about direction, untouched by the magnitude shortfall.

The second check is internal: does the model reproduce the structure the
benchmark reports, depth tracking volumetric energy density while width does not
\cite{Weaver2024}? As a correlation the answer is unstable, over 27 melted
conditions the model gives $r=+0.20$ for width against $+0.77$ for depth, but
width's coefficient runs $+0.68$ to $-0.22$ across five such grids, a property of
the region sampled rather than of the model. The exponents are stable where the
correlation is not. Width and depth are each a power law in $(P,v,d)$ to
$R^{2}\ge0.94$, and a pure function of $\mathrm{VED}=P/(vd^{2})$ must have
exponents along $(1,-1,-2)$: on every grid depth's lie at $\cos 0.64$--$0.74$ of
that ray against width's $-0.11$ to $0.00$, because the two depend on spot size
with opposite signs, the one variable a volumetric density folds away. All 27
points are model output, so this establishes internal consistency, not external
agreement.

\subsection{Laser power and shielding gas are exactly confounded}
\label{sec:confound}

Both faults act \emph{only} on the product (effective absorptivity $\times$
power), in which the Eagar--Tsai temperature rise is linear. Gas attenuates
absorptivity as $A\,e^{-f_{\mathrm{gas}}}$ and power scales as
$P(1+f_{\mathrm{power}})$, so the product is unchanged whenever
$f_{\mathrm{gas}}=\ln(1+f_{\mathrm{power}})$, giving an identical melt pool i.e.
the same temperature at every point, hence identical readings from any instrument
that could be built, the two cancelling bit-identically in double precision on
all three. The
ratio of the two Jacobian columns is $-0.99999$, $-1.00000$, $-0.99993$ across
the sensors, a spread of $7\times10^{-5}$ from finite-difference error on an
exact identity, and the pairwise angle is \textbf{0.00$^\circ$}.

\subsection{The limit is the melt pool, not the sensor count}
\label{sec:melt}

The three observables are not independent. Normalising the Jacobian's sensor Gram
matrix $JJ^{\mathsf T}$ gives pairwise correlations of \textbf{0.99960}
(photodiode--camera), \textbf{0.99508} (photodiode--pyrometer) and
\textbf{0.99477} (camera--pyrometer): a photodiode reading band radiance, a
camera reading pool area in square metres and a two-colour pyrometer reading
kelvin respond to the five faults in patterns agreeing to better than 0.994, their sensitivity vectors lying within $6^\circ$ of one another. Row scalings cancel in the normalisation exactly, so this is the one result here that does
\emph{not} depend on the derived noise figures.

Those correlations set the rank. The whitened Jacobian's singular values are
13.14, 0.0638 and 0.00513: three of the five fault directions are observed, and
the power--gas combination of Section~\ref{sec:confound} lies \textbf{entirely}
inside the unobserved subspace, where power and gas each lie only partly. Of the
three observed, the weakest responds \textbf{2562} times less than the strongest. Adding all four
candidate observables gives seven rows and does not raise the attained rank:
the singular values become $(23.7,\ 9.74,\ 6.0\times10^{-2},\ 5.0\times10^{-6},\
1.4\times10^{-6})$, the fourth at $2\times10^{-7}$ of the leading one, two orders
below the $10^{-6}$ threshold adopted here, so four directions are formally
present but only three are measurements. \textbf{Seven instruments leave power
and gas exactly as confounded as one.}

The reason is structural. Everything the model knows funnels into four physical
channels: absorbed power, scan speed, effective conductivity and preheat,
with three of the five faults crowding onto the first. Any instrument reading
only the melt pool is a function of those four channels. The fault-to-sensor
map cannot exceed rank 4 however many sensors are fitted.

\subsection{Introducing full build context}

With the enlarged observation of Section~\ref{sec:context} the rank rises from 3
to \textbf{5}, the worst pair from $0.00^\circ$ to \textbf{24.90$^\circ$}, power
against gas to $29.34^\circ$, and the condition number falls from 2562 to 50.
All ten pairs clear the $10^\circ$ floor (Table~\ref{tab:matrix}) where none
does per segment.

The mechanism is visible in the form of Equation~\eqref{eq:big}. Each enlarged
observable is a product, $J^{\mathrm{big}}_{(s,j),k} = J_{sk}B_{jk}$: the first
factor is how sensor $s$ responds to fault $k$ at an instant, the second how
strongly fault $k$'s spatial pattern loads onto regressor $j$. The
instantaneous matrix has three rows, so at most three fault directions are
observable whatever the physics. The enlarged matrix has fifteen, but rows
alone buy nothing. What buys separability is that the second factor differs
between faults, so two faults whose sensor signatures are nearly parallel
acquire distinguishable signatures across the regression coefficients.

Three controls establish that this is the mechanism rather than an artefact of
the arithmetic. The first gives all five faults one identical profile. Then $B$
has proportional columns, $B_{jk}=b_j c_k$, and enlarged column $k$ becomes
$c_k\,(J_{:k}\otimes b)$: a scaled Kronecker product of the instantaneous column
with a vector common to every fault. Both the scalar and the common factor
cancel from any angle between two such columns, so the enlarged geometry is the
instantaneous geometry exactly, giving rank 3, condition number 2562, worst pair
$0.00^\circ$, despite the matrix still having fifteen rows. \textbf{The gain is
not dimensionality; it is the profiles differing.}

The second control addresses linearisation. A Jacobian describes infinitesimal
faults, and a real fault has size, so every derivative is replaced by a secant
through the full nonlinear chain. The rank stays 5 at every magnitude from
0.2\,\% to 20\,\% in both signs, but the binding pair ranges $12.82^\circ$ to
$37.95^\circ$ and the condition number 47 to 62: the derivative limit is
therefore \emph{not} a conservative bound, and a $+20\,\%$ fault leaves
$12.82^\circ$, clearing the floor by only 1.3$\times$. The third control
replaces the interpolated response curves with the full model evaluated at all
84 segments, and reproduces the angles to $0.04^\circ$.

\subsection{Robustness}
\label{sec:robust}

The modulation profiles are machine-specific claims rather than measurements, so
the primary result is not an angle but a requirement: how strong a machine's
patterns must be. Scaling every contrast by a common factor, five of the ten
pairs clear the floor at or below 25\,\% of the strength assumed here, and the
hardest, power against speed, at \textbf{45\,\%}. A machine whose patterns are
half as strong as assumed therefore still resolves every pair, and the one
profile that could be calibrated against measurement came back 20\,\%
\emph{stronger} than assumed (Section~\ref{sec:context}).

A second question is which fault follows which build variable. Permuting that
assignment changes the columns of $B$ while leaving each fault's sensor
signature $J_{:k}$ untouched, so the products in Equation~\eqref{eq:big} are
genuinely different vectors and the pairwise angles move. Across all 120
assignments the rank is 5 and all ten pairs clear the floor in every one, while
the worst pair ranges $16.2^\circ$ to $30.1^\circ$. The conclusion is thus
insensitive to the assignment even though the headline angle is not: the method
requires that the patterns differ, not that any particular fault own any
particular one. The headline moves by under half a degree across the full range
of the unmeasured conductivity coupling and across all four noise sets of
Section~\ref{sec:noise}. Segment count is the one hard requirement, at
$11.23^\circ$ for $N=12$ and $22.90^\circ$ for $N=30$ on a near-square plan,
flat thereafter.

That the patterns must differ is a strict requirement, and the same algebra says
why. Give faults $k$ and $l$ identical profiles and $B_{:k}=B_{:l}$, so their
enlarged columns are $J_{:k}\otimes b$ and $J_{:l}\otimes b$ and the angle
between them collapses to the instantaneous one; this is reproduced to better
than $10^{-3}$ degrees for all ten pairs. Every instantaneous angle lies below
the floor, the mildest at $4.63^\circ$, so \textbf{no two faults may share a
spatial signature.} None does on this machine, as the recoater advances once per
layer along a fixed axis while the gas flows along a fixed chamber axis, and the
benchmark measurement of Section~\ref{sec:context} confirms the two alternate at
different scales. Establishing that on another machine is a far weaker
requirement than measuring five contrasts.

Distinctness also tolerates a profile being attributed to the wrong build
variable, which matters because measured layer-thickness variation is reported
to follow part support rather than the recoater cycle assumed here
\cite{Liu2022}. Mixing the geometry fault's support driver into the layer
profile holds layer against geometry above the floor until that profile is
roughly \textbf{84\,\%} support-driven ($12.07^\circ$ at 80\,\%). Mixing the
power fault's build-drift driver never brings power against layer below the
floor at any fraction, because every profile is normalised to unit mean: two
faults sharing a driver give $(1, a z)$ and $(1, b z)$, which are
non-proportional wherever the contrasts $a$ and $b$ differ. Sharing a driver is
not sharing a profile.

\begin{table}[htbp]
\centering
\caption{The diagnosability matrix, in degrees; bold clears the $10^\circ$ floor.
Power/gas reads $0.00^\circ$ in both melt-pool columns, being algebraic. The free
split clears seven pairs and mildly worsens two that stay below the floor either
way; accumulation over $N=84$ clears all ten.}
\label{tab:matrix}
\scriptsize
\begin{tabular}{@{}lrrr@{}}
\toprule
 & \multicolumn{2}{c}{per segment} & over a build \\
\cmidrule(lr){2-3}\cmidrule(lr){4-4}
fault pair & baseline & split camera & baseline \\
\midrule
power / gas      & 0.00 & 0.00 & \textbf{29.34} \\
layer / gas      & 0.32 & \textbf{64.53} & \textbf{44.48} \\
power / layer    & 0.32 & \textbf{64.53} & \textbf{36.26} \\
power / speed    & 0.70 & \textbf{43.43} & \textbf{24.90} \\
speed / gas      & 0.70 & \textbf{43.43} & \textbf{35.87} \\
speed / layer    & 1.03 & \textbf{72.04} & \textbf{46.16} \\
speed / geometry & 3.62 & \textbf{43.28} & \textbf{31.95} \\
power / geometry & 4.31 & 3.08 & \textbf{25.42} \\
gas / geometry   & 4.31 & 3.08 & \textbf{31.39} \\
layer / geometry & 4.63 & \textbf{64.81} & \textbf{40.54} \\
\midrule
pairs clearing the floor & 0/10 & \textbf{7/10} & \textbf{10/10} \\
\bottomrule
\end{tabular}
\end{table}

\subsection{Diagnosability matrix and configuration classes}

Classifying on the enlarged observation alone places all ten pairs in class~1 for
every configuration, a single photodiode included. Both
regimes are therefore reported (Table~\ref{tab:matrix}): \emph{instantaneous},
per-segment decisions for closed-loop control, and \emph{accumulated}, diagnosis
over a build for qualification and drift monitoring. The baseline resolves none of the ten instantaneously but all ten in the
accumulated run, leaving class~3 empty in both.

Five commercial systems were surveyed and present only three distinct observable
sets, reported as classes rather than products because the mapping from product
to observable set rests on secondary literature rather than datasheets: a single
coaxial photodiode (\textbf{A}, rank 1); photodiode plus pool imaging
(\textbf{B}) and two photodiodes in separate bands (\textbf{C}), both rank 2;
the research reference adding pyrometry (\textbf{D}) and C plus upstream power
monitoring (\textbf{E}), both rank 3. \textbf{No configuration built solely from
melt-pool observables separates laser power from gas degradation}, whatever its
sensor count: A through D read $0.00^\circ$ on that pair, for the reason of
Section~\ref{sec:melt}, and none resolves any pair per segment at the derived
noise figures. E is the actionable class, separating the pair at $84.02^\circ$
and resolving four of ten, because a delivered-power measurement sampled
\emph{upstream} of the process zone sees the laser output before the plume
attenuates it.

\subsection{What an added measurement would buy}
\label{sec:added}

Plume and powder-bed sensing cannot be simulated by a model with neither a plume
nor a powder layer, so each candidate is represented as an idealised unit row on
its cause. One such row raises the rank to 4 and clears the four pairs involving
that cause. \textbf{Build context outperforms two idealised, perfect direct
measurements of the causes themselves}: gas and layer sensors together give rank
5 but $0.70^\circ$ and seven of ten, against $24.90^\circ$ and ten of ten at no
cost. Among melt-pool observables the most valuable addition is free, as
reporting the camera's two axes separately rather than as an area is worth
\textbf{+23.89$^\circ$} on the pair that binds the baseline and is the only
change resolving anything \emph{per segment}, taking that count from none of ten
to \textbf{seven of ten} (Table~\ref{tab:matrix}), while a short-wave infrared
photodiode and a wide-band pyrometer are worth nothing. Width is the least
correlated of the candidates with the coaxial photodiode, 0.779 against 0.979 for
length, and an area averages it away; \emph{a sensor added to break a confound
must observe upstream of the shared channel, where a sensor reading the melt pool
adds a row, not a direction.}

%% ===========================================================================
\section{Discussion}
%% ===========================================================================

\subsection{Relation to reported monitoring performance}

Comparing predicted limits against reported accuracies proves largely
unrealisable, and the reason is itself a result. Of 19 works surveyed, covering in-situ
LPBF monitoring with a stated input and output over 2019--2025, seven classify
outcomes and only \textbf{three estimate a process input from an in-situ
signal}, one of those a lumped energy-density class; none analyses whether two
causes are separable. That is not an accident of what the field chose to study.
Outcomes are functions of the melt pool, and the melt pool is a function of the
four channels of Section~\ref{sec:melt}, so outcome classification cannot
separate two causes sharing a channel, for the same reason no melt-pool
instrument can. No published result contradicts the predicted limits, and
state-of-the-art parameter estimation predicts precisely the lumped group this
analysis identifies as what survives confounding \cite{kNN2025}.

The one external check available on this model's sensor response turns on an
assumption worth stating before the result. The plume here is a \emph{fixed}
optical depth, independent of power and speed and emitting nothing of its own,
so the model admits no mechanism by which its scan-speed response could
strengthen with laser power. Reijonen measures that response on a commercial
coaxial photodiode as $y=ax^{-b}$, reporting $b=0.337$ at 400\,W over
400--1200\,mm\,s$^{-1}$ and the exponent growing with power, ``almost
non-existent'' at 100\,W \cite{Reijonen2024}. Fitting the same form to this
model over their range, which is the comparable quantity rather than the Jacobian
entry $-0.483$, gives \textbf{0.468} at 400\,W, some 39\,\% too steep, and
0.51, 0.44 and 0.47 at 100, 285 and 400\,W: flat, as the assumption requires.
Reijonen attribute the trend to the vapour plume, absent in conduction mode at
low power, which is the mechanism omitted here. \textbf{That omission is the
same one that makes the power/gas degeneracy exact} (Section~\ref{sec:confound}):
their measurements show the real attenuation depending on both parameters and
even reversing sign with scan speed at low gas flow, so the reported
$0.00^\circ$ is an upper bound on the instantaneous confounding on external
evidence rather than by assertion.

The comparison crosses material (316L against IN718), detection band (a
shortpass below 800\,nm against 700--1000\,nm) and bed condition
(30\,\textmu m powder layers against bare plate). In the model's defence it is
calibrated at 285\,W, where all seven benchmark cases keyhole, so the fixed
optical depth is a local approximation \emph{inside} the regime where a plume
exists rather than a claim about conduction-mode melting. On the power response,
Reijonen report a linear fit $y=ax+b$ without quoting the intercept; relative
sensitivity is then $aP/(aP+b)$, equal to one only through the origin, so this
model's $+1.52$ cannot be called a conflict until that intercept is known.

\subsection{Limitations}
\label{sec:limits}

The thermal model is a conduction solution with an analytic keyhole patch: it
only partly reproduces the Guo validation, its fitted parameters sit at their bounds, and
the absorptivity, conductivity and heat capacity correlations of 0.93--0.96
reported during Bayesian calibration of the same model family \cite{Hon2021} are
the pathology this paper studies appearing in its own calibration. Three
modelling choices then make the instantaneous case \emph{more} confounded than
reality: the gas fault as pure optical depth omits a convective channel; treating
plume attenuation as independent of power and speed is what makes the degeneracy
exact; and modelling upstream power monitoring as a melt-pool-derived signal is
likewise pessimistic.

\textbf{The $0.00^\circ$ is therefore an upper bound, and its exactness is an
artefact of the one-channel gas model.} Gas here acts only on absorbed power,
precisely the channel a power fault acts on, which is what makes the two
algebraically identical. A real plume also scatters: Reijonen et al.\ report
reduced gas flow giving loss of penetration together with \emph{increased} melt
pool width \cite{Reijonen2020}, where attenuation alone gives a decrease of
$-10$\,\% here at $f_{\mathrm{gas}}=0.4$. Adding that broadening as a spot-size
channel, which power does not possess, opens the pair to \textbf{2--3$^\circ$}
at the strength needed to reverse the width sign, and leaves the suite's worst
pair unmoved at $0.32^\circ$. The exactness is a property of the model; the
inseparability is not. The rank limit is unaffected either way, deriving from
channel counting rather than the form of this coupling.

The geometry inversion fails, returning $+687$\,K of preheat because a pad
differs from a single track by remelting as well as preheat, and no dataset yet
offers a build with \emph{deliberately induced} machine faults: every inversion
here reads a healthy build and extrapolates, which is the natural follow-on
experiment. The noise derivation rests on one assumption, that every radiometric
channel carries the same fractional noise, bounded by the gain-dominated sweep.
The five fault magnitudes are physical but not on a common severity scale, so
comparisons \emph{across pairs at fixed magnitude} are not meaningful, and all
results are diagnosability statements, meaning what a suite can distinguish in
principle, at a given relative noise, under a given forward model, not that
any algorithm achieves it.

%% ===========================================================================
\section{Conclusions}
%% ===========================================================================

Laser power deviation and shielding gas degradation produce identical melt pools
in this model, confounded at $0.00^\circ$ by algebra rather than noise, so no
instrument reading the melt pool separates them from a single instant. The
binding limit is the melt pool's four physical channels, not the instrumentation:
at noise derived from each sensor's transduction chain, seven observables attain
rank 3 against that bound of 4 and leave the worst pair at $0.00^\circ$. Deriving
those figures rather than assuming them matters, whitening being half of the
measurement geometry.

Build structure already present dissolves the limit at no hardware cost, taking
rank 3 $\rightarrow$ 5, worst pair $0.00^\circ \rightarrow 24.90^\circ$ and all ten pairs
resolved where none is per segment, surviving at half the assumed profile
strength, across all 120 permutations, under all four noise sets and at fault
magnitudes from 0.2\,\% to 20\,\%, with a shared-profile control reproducing the
instantaneous result exactly. It outperforms two idealised direct measurements of
the causes themselves and requires only that no two faults share a spatial
signature; those patterns being machine-specific, the deliverable is the method
plus the strength a machine needs, with the inversion of
Section~\ref{sec:context} as the recipe for calibrating it.

For real-time control an upstream measurement of delivered laser power is the
recommendation, separating power from gas at $84^\circ$ where no melt-pool
configuration separates them at all; among melt-pool observables the most
valuable addition is free, reporting the camera's two axes separately being worth
$+23.89^\circ$ and the only change resolving anything per segment. Post-build
diagnosis needs no new hardware, real-time isolation does, and crude sensing of
the right cause suffices. The method needs only a forward model and a sensor
description, and answers at specification time what is otherwise settled after
the hardware is bought.

%% ===========================================================================
\section*{Declaration of generative AI and AI-assisted technologies in the
manuscript preparation process}

During the preparation of this work the author used Claude Code (Anthropic) in
order to carry out part of the formal analysis and to edit the original draft. After using
this tool, the author reviewed and edited the content as needed and takes full
responsibility for the content of the published article.
%% ===========================================================================
%% TODO before submission: add funding and acknowledgements section if required.

%% ===========================================================================
\begin{thebibliography}{99}
\scriptsize
\setlength{\itemsep}{-2pt}\setlength{\parskip}{0pt}

\bibitem{Krysander2008}
M. Krysander, E. Frisk, Sensor placement for fault diagnosis, IEEE Trans. Syst. Man Cybern. A 38 (6) (2008)
1398--1410.

\bibitem{Scime2019}
L. Scime, J. Beuth, Using machine learning to identify in-situ melt pool
signatures indicative of flaw formation in a laser powder bed fusion additive
manufacturing process, Addit. Manuf. 25 (2019) 151--165.

\bibitem{Raj2023}
A. Raj et al., Modeling spatial variations in co-axial melt pool monitoring signals in laser
powder bed fusion, J. Manuf. Process. 89 (2023) 24--38.

\bibitem{Blanc2024}
C. Blanc, A. Ahar, K. De Grave, Reference dataset and benchmark for
reconstructing laser parameters from on-axis video in powder bed fusion of bulk
stainless steel, arXiv preprint arXiv:2409.12475 (2024).

\bibitem{Reijonen2024}
J. Reijonen, Effect of laser parameters and shielding gas flow on co-axial
photodiode-based melt pool monitoring signals in laser powder bed fusion,
Addit. Manuf. Lett. 11 (2024) 100244.

\bibitem{Reijonen2020}
J. Reijonen et al., On the effect of shielding gas flow on porosity and melt
pool geometry in laser powder bed fusion additive manufacturing, Addit. Manuf.
32 (2020) 101030.

\bibitem{Liu2022}
T. Liu et al., In-situ infrared thermographic inspection for local powder layer
thickness measurement in laser powder bed fusion, Addit. Manuf. 55
(2022) 102873.

\bibitem{Bugatti2022}
M. Bugatti, B. M. Colosimo, The intelligent recoater: A new solution for in-situ
monitoring of geometric and surface defects in powder bed fusion, Addit. Manuf. Lett. 3 (2022) 100048.

\bibitem{kNN2025}
Estimation of laser powder bed fusion process parameters from melt pool images
using k-nearest neighbours, Int. J. Adv. Manuf. Technol. (2025). doi:10.1007/s00170-025-15591-y.

\bibitem{Guo2019}
Q. Guo et al., In-situ
characterization and quantification of melt pool variation under constant input
energy density in laser powder bed fusion additive manufacturing process,
Addit. Manuf. 28 (2019) 600--609.

\bibitem{Weaver2024}
J. S. Weaver et al., Cross-sectional melt pool geometry of laser scanned tracks and pads on nickel
alloy 718 for the 2022 additive manufacturing benchmark challenges, Integr. Mater. Manuf. Innov. 13 (2024) 363--379.

\bibitem{Hon2021}
P. Honarmandi et al., A rigorous test and improvement of the
Eagar-Tsai model for melt pool characteristics in laser powder bed fusion
additive manufacturing, Addit. Manuf. 47 (2021) 102300.

\bibitem{Trapp2017}
J. Trapp et al., In situ absorptivity
measurements of metallic powders during laser powder-bed fusion additive
manufacturing, Appl. Mater. Today 9 (2017).
doi:10.1016/j.apmt.2017.08.006.

\end{thebibliography}

\end{document}
